% Reserved C.3br. Complete draft; unregistered pending independent audits.
\paragraph{C.3br. Component families for a triple-free exceptional surface
(complete draft).}
Let $k$ be algebraically closed of characteristic $p>0$ and let
$E=\bigcup_{i=1}^cE_i$ be a projective SNC surface with smooth connected
components. Assume that there are no triple intersections.
Index every connected double curve separately as $C_e$, and orient its
edge $e=(i,j)$. Put
\begin{align}
 \rho:\bigoplus_i\operatorname{Pic}(E_i)_{\mathbf Q}
   &\longrightarrow\bigoplus_e\operatorname{Pic}(C_e)_{\mathbf Q},
 & (L_i)_i&\longmapsto
       (L_i|_{C_e}-L_j|_{C_e})_{e=(i,j)},                 \tag{TF.1}\\
 S(E)&=F_\gamma^2K_0(E)_{\mathbf Q}. \notag
\end{align}
As in EP, $J_K(E)$ denotes the subgroup generated by endpoint
differences of actual classes in $F_\gamma^2K_0(E\times P)_{\mathbf Q}$
with $P$ a smooth connected projective curve.

If $\rho$ is surjective, then the component Chern-character map is an
isomorphism
\begin{equation}
 \varphi:S(E)\xrightarrow{\ \sim\ }
                     \bigoplus_i\operatorname{CH}_0(E_i)_{\mathbf Q},
 \qquad
 \varphi(J_K(E))=\bigoplus_i\operatorname{CH}_0(E_i)^0_{\mathbf Q}.
                                                               \tag{TF.2}
\end{equation}
Consequently,
\begin{equation}
                       S(E)/J_K(E)\simeq\mathbf Q^c.     \tag{TF.3}
\end{equation}
In particular, for two components $E=S\cup_C T$, the hypothesis
holds if the sum of the restrictions
\begin{equation}
 \operatorname{Pic}^0(S)_{\mathrm{red}}\times
 \operatorname{Pic}^0(T)_{\mathrm{red}}\longrightarrow
                          \operatorname{Pic}^0(C)        \tag{TF.4}
\end{equation}
is surjective. It suffices that either factor already surject.
It also holds whenever $C\simeq\mathbf P^1$.

Now assume the exact threefold resolution hypotheses of C.3bm:
$R$ is the strict henselization at a closed $k$-point of a normal
local ring essentially of finite type over $k$, with $\dim R=3$ and
regular puncture;
$Z\to\operatorname{Spec}R$ is projective, its source is regular of pure
dimension three, and it is an isomorphism off the closed point; $E$ is its full
reduced closed fibre. Suppose that
\[
 D=\sum_i a_iE_i,\qquad a_i>0,\qquad
                         \mathcal O_E(-D)\text{ is ample}.
\]
If $\rho$ is surjective, then
\begin{equation}
                    B^2(Z\setminus E)_{\mathbf Q}=0.     \tag{TF.5}
\end{equation}
In this resolved setting, surjectivity of $\rho$ is automatic if every
double curve is rational. Thus (TF.5) holds for every such triple-free
rational-double-curve configuration, with no condition on the
Albanese kernels or geometric genera of its components.
The assertion supplies no resolution-existence theorem and does not
include configurations with triple points.

\begin{proof}
\emph{1. The exact two-column descent sequence.}
For every actual smooth projective parameter scheme $P$, closed-union
descent gives the homotopy fibre sequence
\begin{equation}
 K(E\times P)_{\mathbf Q}\longrightarrow
       \bigoplus_iK(E_i\times P)_{\mathbf Q}
   \xrightarrow{\ r\ }\bigoplus_eK(C_e\times P)_{\mathbf Q}.
                                                               \tag{TF.6}
\end{equation}
Its maps are alternating derived restrictions. There is no triple
term or additional coherent nullhomotopy to specify.
For completeness, each affine binary closed union is a Milnor square.
Theorem C of \citet[printed p.~30]{rootGeisserHesselholt2011} makes its
bi-relative groups zero after $p$-inversion in every integer degree;
Theorem A permits any nilpotent reductions in the induction.
Zariski descent \citep[Theorem 8.1 and proof, p.~367]
{D_ThomasonTrobaugh1990} globalizes this assertion. Finite induction
over the components gives (TF.6), naturally in $P$. Distinct double
curves on the same component are disjoint because triple points
are excluded.

The projective schemes $E$ and $E\times P$ have finite gamma
filtrations in degrees at most two and three, respectively.
Their rational Adams eigenspace projectors therefore exist
\citep[Chapter II, Proposition 4.9, Theorem 4.10, and
Remark 8.8.5]{D_Weibel2013KbookII}. We use these projectors only
in $K_0$, not an unproved splitting of a fibre spectrum.
They commute with all the restriction maps in (TF.6).

\emph{2. Surjectivity of the component map.}
On a smooth projective surface, the pure weight-two Chern character
identifies $F_\gamma^2K_0(E_i)_{\mathbf Q}$ with
$\operatorname{CH}_0(E_i)_{\mathbf Q}$. On a smooth curve,
$F_\gamma^2K_0(C_e)=0$. A tuple of pure weight-two component classes
therefore restricts to zero in the last group of (TF.6) for a point
parameter. Exactness supplies a class in $K_0(E)_{\mathbf Q}$.
Applying its weight-two projector preserves the tuple and proves
surjectivity of $\varphi$.

\emph{3. The full edge kernel is zero under (TF.1).}
For a smooth connected proper curve $C$ over $k$, localization and
the divisor valuation map give
\[
 \bigoplus_{x\in C(k)}K_1(k)\longrightarrow K_1(C)
 \longrightarrow k(C)^*\longrightarrow
                       \bigoplus_{x\in C(k)}\mathbf Z .
\]
The image in $k(C)^*$ is $k^*$, since a rational function with no
zero or pole on a proper connected curve is constant. This image
is split by the constant units. The remaining subgroup is generated
by the point pushforwards
\[
          u[\mathcal O_x]=u\bigl([\mathcal O_C(x)]-1\bigr),
                         \qquad u\in k^*.
                                                               \tag{TF.7}
\]
The equality holds in $K_1(C)$: on a curve the two rank-zero
line-bundle descriptions of the point class agree, and the
projection formula applies to a constant unit.
These are exactly the weight-two terms after rationalization;
constant units have weight one.
The precise localization and projection inputs are
\citet[Section 7, Propositions 2.10 and 3.2, equation (3.3),
original pp.~126--127]{rootQuillen1973}.
The same Picard-times-units generation is stated in the proof of
\citet[Lemma 5.2, pp.~395--396]{D_PedriniWeibel1994}.

Let $s\in S(E)$ have zero component image. The long exact sequence
of (TF.6) expresses it as the boundary of an edge $K_1$ vector.
Applying $(\psi^2-2)/2$ to that vector does not change its boundary,
because $\psi^2s=4s$, and kills its constant-unit part.
Here Adams operations commute with this boundary by the
arbitrary-scheme iterated-homotopy-fibre and delooping construction
in \citet[p.~258, paragraph preceding Lemma 2.1]{frC3bcLevine1994}.
We use that functorial construction and its agreement with the
classical operations on $K_0$; we do not apply the regular-ambient
hypothesis of Lemma 2.1 to the singular scheme $E$.
It is consequently a sum of vectors of the form (TF.7).
For each fixed unit in such a finite sum, surjectivity of $\rho$
lifts the edge Picard tuple to rational component Picard classes.
Represent the latter by finite rational combinations of line bundles
and multiply their rank-zero $K_0$ classes by that unit.
Their restrictions equal the required vector, since
$L\mapsto[L]-1$ is additive on the Picard group of a smooth curve.
They are actual component $K_1$ classes, whose boundary is zero
in (TF.6). Hence $s=0$ and $\varphi$ is an isomorphism.
No surjectivity of a component Picard map on an entire family
scheme is asserted here.

\emph{4. Move an arbitrary component difference by a fixed proper curve.}
Write $\partial E_i=\bigcup_{e\ni i}C_e$.
Every zero-cycle on $E_i$ is rationally equivalent to one supported
away from $\partial E_i$. Here is a direct moving argument.
Choose a smooth projective curve through any unwanted point,
not contained in $\partial E_i$. Its intersection with $\partial E_i$
is finite. On that curve, choose a point outside this finite set and
a sufficiently large multiple of it. Riemann--Roch then supplies
a section of the corresponding twist of the unwanted point divisor
which is nonzero at every point of the finite set. Its zero divisor
is supported in the complement, giving the required rational
equivalence. Such smooth curves exist by a sufficiently high
very ample system with the one prescribed point.

It therefore suffices to treat $[y]-[x]$ with
$x,y\in E_i\setminus\partial E_i$. Choose a smooth connected
projective curve $P\subset E_i$ through $x,y$, transverse to all
double curves. To justify the choice in every characteristic, use
a sufficiently high very ample system separating the required values
and first jets. After the two value conditions, singularity at any
other point imposes three independent conditions over a surface,
while tangency to a fixed smooth double curve imposes two over a
curve. The closures of these bad incidence images are proper by
the dimension count. At $x,y$ the free first jets prevent singularity.
Connectedness follows from Serre duality and Serre vanishing for
the high ample divisor. This is the same elementary Bertini argument
as \citet[Lemma 33.47.3]{D_StacksBertiniUniformHigher}.
Write $o,t\in P(k)$ for its parameter points $x,y$.

The graph of $P\hookrightarrow E_i$ is a regular codimension-two
closed subscheme of $E_i\times P$. Take its perfect $K_0$ class and
its pure weight-two projection; set all other component classes
to zero. Call the resulting component tuple $\alpha$.
Its two component fibres are $[x]$ and $[y]$ on $E_i$ and zero
on every other component.

The graph meets each $C_e\times P$ in finitely many transverse
points $(c_{e,\ell},p_{e,\ell})$. Consequently its alternating
restriction to the edge products is
\begin{equation}
 (r\alpha)_e=
    \sum_\ell\varepsilon_{i,e}
           [\mathcal O_{c_{e,\ell}}]\boxtimes
                              [\mathcal O_{p_{e,\ell}}],
             \qquad \varepsilon_{i,e}\in\{1,-1\}.       \tag{TF.8}
\end{equation}
Neither endpoint $o,t$ occurs among the $p_{e,\ell}$.
This is an equality in perfect $K_0(C_e\times P)_{\mathbf Q}$:
the intersection is Tor-independent and its zero-dimensional
class is pure weight two. It is not merely an equality in a
coherent-sheaf quotient of the singular E.

For each of the finitely many parameter values $p$ in (TF.8), use
$[\mathcal O_c]=[\mathcal O_{C_e}(c)]-1$ to regard its edge
coefficients as a tuple in $\bigoplus_e\operatorname{Pic}(C_e)_{\mathbf Q}$.
Surjectivity of $\rho$ gives rational component divisor classes
$L_{j,p}$ with exactly that restriction tuple.
Let $\ell_{j,p}\in K_0(E_j)_{\mathbf Q}^{(1)}$ be their pure
weight-one classes. Their edge restrictions are the corresponding
rank-zero line classes. Form the actual component families
\begin{equation}
 v_j=\sum_p\ell_{j,p}\boxtimes[\mathcal O_p]
            \in K_0(E_j\times P)_{\mathbf Q}^{(2)}.      \tag{TF.9}
\end{equation}
They have zero fibres at both $o$ and $t$, and $rv=r\alpha$.
Thus the corrected component tuple $\alpha-v$ lies in the kernel
of the last map in (TF.6). It has an actual lift to
$K_0(E\times P)_{\mathbf Q}$, whose weight-two projection belongs
to $F_\gamma^2K_0(E\times P)_{\mathbf Q}$ and has the same component
restrictions. Its endpoint difference belongs to $J_K(E)$ and maps
under $\varphi$ to $[y]-[x]$ on the selected component and zero
on all the others.

This construction uses only finitely many rational divisor lifts
at fixed parameter values. It does not replace the fixed product
by a varying neighbourhood, and it does not demand an actual
Poincare lift for every component Picard family.

\emph{5. Identify exactly the relation group.}
Step~4 proves that $\varphi(J_K(E))$ contains every component
degree-zero zero-cycle. Conversely the component degrees are
constant in an actual proper perfect family. For example, restrict
its class to $E_i\times P$, take $\operatorname{ch}_2$, and push
forward to $P$; the resulting codimension-zero Chow class is a
constant multiple of $[P]$, and refined restriction gives that same
degree on each fibre. Thus
$\varphi(J_K(E))=\bigoplus_i\operatorname{CH}_0(E_i)^0_{\mathbf Q}$.
Together with Step~3 this proves (TF.2)--(TF.3), including arbitrary
component Albanese kernels.

For two components, surjectivity in (TF.4) supplies all degree-zero
Picard classes on C. An ample line bundle from either component has
positive degree on C, so it supplies the remaining degree coordinate
after tensoring with $\mathbf Q$. This proves the two-component
criterion. For $C\simeq\mathbf P^1$, there is no degree-zero Picard
coordinate, so no additional hypothesis is needed.

\emph{6. The negative exceptional geometry forces the rational-edge
criterion.}
Assume now the resolved setting and that all $C_e$ are rational.
For $e=(i,j)$ put
\[
 n_{i,e}=\deg N_{C_e/E_i}=C_e^2|_{E_i},\qquad
 n_{j,e}=\deg N_{C_e/E_j}=C_e^2|_{E_j}.
\]
The two normal bundles are interchanged with the normal bundles
of the exceptional components inside Z. Because no other
component meets $C_e$, one has
\begin{equation}
              \deg(D|_{C_e})=a_i n_{j,e}+a_j n_{i,e}<0.
                                                               \tag{TF.10}
\end{equation}
The strict inequality follows from ampleness of $-D$.
At least one of $n_{i,e},n_{j,e}$ is therefore nonzero.
If $n_{i,e}\ne0$, the rational divisor
$n_{i,e}^{-1}C_e$ on $E_i$ restricts with degree one to that edge;
it restricts trivially to every other edge on $E_i$, since those
curves are disjoint. If the nonzero self-intersection is on
$E_j$, use $-n_{j,e}^{-1}C_e$ there to account for the orientation.
Since $\operatorname{Pic}(C_e)_{\mathbf Q}\simeq\mathbf Q$ by
degree, these give an explicit right inverse to (TF.1).
The argument treats multiple disjoint edges separately and places
no restriction on cycles in the dual graph.

\emph{7. Pass to the puncture quotient.}
C.3bm gives
\[
 B^2(Z\setminus E)_{\mathbf Q}
   =S(E)/(J_K(E)+\theta\operatorname{CH}_1(E)_{\mathbf Q}).
\]
By (TF.3), only the component-degree coordinates remain.
Choose an ample H on E and let $M_{ij}$ be the degree on $E_i$
of the exceptional restriction of the curve class $H_j$ on $E_j$.
The matrix is symmetric, its off-diagonal entries are nonnegative,
and $Ma$ is strictly negative because
$(Ma)_i=\deg(D|_{E_i}\cdot H_i)<0$.
For $A=\operatorname{diag}(a_i)$, the identity
\[
 x^{\mathsf T}AMA x
  =\sum_i\left(\sum_j(AMA)_{ij}\right)x_i^2
       -\sum_{i<j}(AMA)_{ij}(x_i-x_j)^2
\]
is strictly negative for nonzero real x. Thus M is invertible
over $\mathbf Q$, and exceptional curve classes span all of
$\mathbf Q^c$. This proves (TF.5).
\end{proof}

\emph{Exact remaining scope.}
The intrinsic criterion removes all component degree-zero classes,
not only their Albanese images, but requires the full Picard
restriction surjectivity and the absence of triple strata.
Negative rational double curves give that surjectivity explicitly.
For positive-genus double curves it is a separate hypothesis.
With triple points, pairwise equality of edge classes alone need
not give an actual global perfect class; no such extension is
included here. The higher-dimensional analogue is also not part
of this statement.
